%% ============================================================
%% preamble/verse.tex — parallel-text verse (module: verse)
%%
%% Purpose      : facing original/translation verse in two synced
%%                columns (70/30 by default):
%%                  \begin{parallelverse}
%%                    \begin{verse} original lines \end{verse}
%%                    \switchcolumn
%%                    \begin{verse} translation \end{verse}
%%                    \stanzasync
%%                    ... next stanza pair ...
%%                  \end{parallelverse}
%%                plus per-5 verse line numbers and hanging indents.
%% Dependencies : paracol + verse (loaded here), colors.tex (bk-folio
%%                for line-number ink). Loaded only when
%%                \ifBookModuleVerse is true.
%% Provenance   : iliad-brainrot / ovid-brainrot parallel-text setup
%%                (docs/research/iliad-brainrot.md), incl. the
%%                \setchapterheads markboth workaround.
%%
%% ENGINE NOTE  : this module is engine-safe. Greek/Latin language
%%                support (polyglossia) exists ONLY under XeLaTeX —
%%                book.yaml enforces engine: xelatex when
%%                modules.verse is on (config-schema.md) — but the
%%                guard below keeps the module compilable under
%%                LuaLaTeX too, so flipping the toggle never produces
%%                a mystery "polyglossia wants XeTeX" abort.
%% ============================================================

\RequirePackage{iftex}   % \ifxetex below; explicit — don't rely on callers
\usepackage{paracol}
\usepackage{verse}

% Original text gets the wide column; translation reads narrower.
\columnratio{0.70}

% polyglossia (ancient Greek shaping, Latin hyphenation) is XeTeX-only
% territory in this stack; under LuaLaTeX the module still typesets,
% just without the extra language machinery.
\ifxetex
  \usepackage{polyglossia}
  \setdefaultlanguage{english}
  % Books add e.g.: \setotherlanguage[variant=ancient]{greek}
  %                 \newfontfamily\greekfont{...}[Script=Greek]
\fi

% Verse pages must NOT flush to a common bottom edge — stanza glue is
% meaningful. This overrides styling.tex's prose default (documented
% policy: latex-pipeline.md §5).
\raggedbottom

% ----------------------------------------------------------------------
% parallelverse — a paracol pair. \stanzasync ends a stanza pair and
% forces both columns to realign before the next pair (a plain
% \switchcolumn drifts once stanza lengths diverge).
% ----------------------------------------------------------------------
\newenvironment{parallelverse}
  {\begin{paracol}{2}}
  {\end{paracol}}
\newcommand{\stanzasync}{\switchcolumn*}

% Line numbers every 5 lines, in the margin ink used for folios; run-in
% verse lines hang by 2em so numbered lookups don't misparse wraps.
\poemlines{5}
\verselinenumfont{\footnotesize\color{bk-folio}}
\setlength{\vindent}{2em}

% ----------------------------------------------------------------------
% \setchapterheads{verso}{recto} — iliad/ovid workaround: paracol
% swallows the \markboth that \chapter issues when a chapter starts
% inside/near a column context, leaving stale running heads. Verse
% chapters call this right after \chapter to set marks explicitly.
% ----------------------------------------------------------------------
\newcommand{\setchapterheads}[2]{\markboth{#1}{#2}}
